\paragraph{Performance measurement scope.}
The performance tables use an H100 PCIe, PyTorch~2.14, Triton~3.8 and
\texttt{mamba\_ssm}~2.3.2.post1. Main timings and memory use bf16;
the separately labelled appendix ablation uses fp32 input/output.
The Mamba recoveries use copied official weights and pre-timing output checks;
the Mamba~3 recovery matches the installed baseline's softplus activation.
The exact kernel-scope boundaries and gradient leaves are specified per pair
in the accompanying harness method; mixer measurements cover the full module.
Table~\ref{tab:bench_cost} rows implement different functions and are not recovery speed ratios.
The rank rows use model head width $P=64$, with real kernel widths $G=32$
at R2 and $G=16$ at R4; these are not fixed-$G$ rank scaling experiments.

\paragraph{Selection and aggregation.}
Each comparison group is measured for three rounds with alternating arm order.
The reported time is the median of the three round values; a forward-and-backward
round is an event-timed mean of ten calls after eight warm-up calls. Selection
uses the first attempt passing the validity checks and a per-group canary spread
of at most $1.15$. Selection does not use the measured latency. Whole-run drift
is recorded rather than gated. The Table~\ref{tab:app_parity_summary} summaries use geometric means and
min--max ranges over four Large or six LM shapes; HM results beyond the
five main shapes retain the earlier buffer schedule, as labelled. Noise-only failures with no
accepted attempt remain explicitly n/a.

\paragraph{Memory.}
Extra saved is retained allocator growth between a gradient-enabled
forward and its backward, excluding parameters, caller-held inputs and the
output. The harness checks release after backward and freeing leaf gradients.
Peak is the maximum allocated memory during one forward and backward above
the resident baseline of parameters, caller inputs and output cotangent,
and includes outputs, temporary workspace and newly allocated gradient buffers.

Tables~\ref{tab:bench_parity} and~\ref{tab:bench_cost} additionally report
unique nonparameter backing storage held for backward, measured in separate
memory-only runs on the same implementation. This retained total includes
caller inputs and outputs when saved, deduplicates aliases, and includes
custom context tensors and inputs held for regeneration without unpacking
saved tensors. It excludes resident parameters. Allocation and post-backward
release checks validate the accounting. HM and paired Mamba~1 timings and
peaks use a complete five-shape campaign; retained totals match the earlier
memory-only measurements exactly. SSD-MIMO and Mamba~2/3 retain their accepted
measurements; the memory-only runs reproduce their matched-pair peaks within $0.5\%$.
The remaining appendix tables retain the extra-saved metric. Retained storage
and peak allocation overlap and must not be added; neither is total-model
training memory.

\paragraph{Rank configurations and coverage.}
The real R2O2 and R4O4 configurations use forward/backward chunks $32/32$
and $16/16$, respectively. Chunk selection uses production-mixer workloads
at $N=128$, over 18 shared shapes for R2 and six LM shapes for R4.
Launch configurations are selected using five rounds on two representative
workloads per rank at $N=64$, then measured at the paper geometry $N=128$.
This is a targeted configuration screen, not an exhaustive search for an
optimum on every measured shape. Tuning and compilation precede timing.
Complex R1/R2 retain their measured configurations. Complex R4 uses training
chunks $16/16$ at kernel width $16$. Its launch configurations use the generic
light search for this chunk size during warmup, outside the reported timings.
The real R2/R4 and complex R4 training configurations have kernel and eager-mixer
forward-and-backward measurements on all five paper shapes, and saved/peak
memory at the LM shape; other training entries are n/a. Complex R4's
forward-only results retain the original measurements because its inference
configuration is unchanged.

\paragraph{Complex R4 speed and memory.}
A bounded configuration screen compared chunks $16$ and $32$; a separate run
supplies the reported table entries. The same-weight fp64 output/input-gradient
check and gradient-presence check passed. At the LM shape, chunk $16$ reduces
kernel forward-and-backward time from $70.89$ to $63.71$\,ms, while kernel
peak memory rises from $2549$ to $3125$\,MiB and mixer peak memory from
$2818$ to $3394$\,MiB. Extra-saved memory is unchanged ($974$ and
$1359$\,MiB). This is a speed/peak-memory tradeoff, not a uniform improvement.

\paragraph{Limits of the evidence.}
The Mamba~3 recovery's LM kernel forward-and-backward cell varied by about
$10\%$ between earlier processes (approximately $2.90$ versus $3.19$\,ms), while
the official arm was stable; the cause remains unresolved. HM measurement records
include launch configurations but some calibration-source fields are null, so
complete calibration-source provenance is not claimed for them. The real rank rows additionally record calibration file hashes, active groups and cached
launch configurations. Their short numerical check compares outputs and input
gradients against the same-weight fp64 Torch reference; it is not an exhaustive
numerical-validation campaign.

\paragraph{HM temporary-buffer schedule.}
The optimized backward schedule reduces redundant gradient-buffer slots and
releases scratch before subsequent adjoint operations, reducing overlapping
lifetimes. It does not reduce the activations or checkpoints retained for
backward. Table~\ref{tab:hm_workspace_tradeoff} compares both schedules on all
five main shapes using absolute measurements from each. The campaign includes
forward-only, eager and CUDA-graph training times, kernel/mixer peak allocation,
and retained storage. This buffer-schedule comparison is separate from the copied-weight fp64
recovery studies of Appendix~\ref{app:numerics_exact}.
The additional five LM shapes and the small/medium HM workloads were not
remeasured; their earlier-schedule results remain explicitly labelled.
